\chapter{Probability and Statistics for Machine Learning}
\label{ch:probability}

\epigraph{``All models are wrong, but some are useful.''}{George E. P. Box~\cite{box1976science}}

\learninggoal{
	\item Review the core probability concepts: random variables, distributions, expectation, Bayes' rule.
	\item Understand maximum likelihood and maximum a posteriori estimation.
	\item Distinguish between parametric and non-parametric methods.
	\item Apply hypothesis testing and confidence intervals to model comparison.
}

\section{Random Variables and Distributions}

A \emph{random variable} $X$ is a variable whose value is determined by a random process. The \emph{probability distribution} describes how likely each value is.

\begin{definition}[Probability Distribution]
	A \textbf{probability distribution} assigns a probability to each possible outcome of a random variable. For discrete $X$, we use a probability mass function $p(x) = P(X = x)$. For continuous $X$, we use a probability density function $f(x)$ where $P(a \leq X \leq b) = \int_a^b f(x)\,dx$.
\end{definition}

\subsection{Key Distributions for ML}

\begin{longtable}{p{3cm}p{4cm}p{5cm}}
	\toprule
	\textbf{Distribution}     & \textbf{Parameters}             & \textbf{Uses in ML}               \\
	\midrule
	Bernoulli$(p)$            & $p \in [0,1]$                   & Binary classification (coin flip) \\
	Binomial$(n, p)$          & $n$ trials, $p$ success         & Count of successes                \\
	Normal$\N(\mu, \sigma^2)$ & mean $\mu$, variance $\sigma^2$ & Noise model, continuous targets   \\
	Exponential$(\lambda)$    & rate $\lambda$                  & Waiting times                     \\
	Poisson$(\lambda)$        & rate $\lambda$                  & Count data                        \\
	Beta$(\alpha, \beta)$     & shape $\alpha, \beta$           & Prior for Bernoulli/Binomial      \\
	Gamma$(k, \theta)$        & shape $k$, scale $\theta$       & Prior for precision               \\
	\bottomrule
\end{longtable}

The normal distribution deserves special attention. Its density is:

\[
	p(x) = \frac{1}{\sqrt{2\pi\sigma^2}} \exp\left(-\frac{(x - \mu)^2}{2\sigma^2}\right)
\]

\subsection{Expectation and Variance}

\begin{definition}[Expectation]
	The \textbf{expected value} of a function $g$ of random variable $X$ is:
	\[
		\E[g(X)] = \begin{cases}
			\sum_x g(x) p(x)     & \text{if } X \text{ is discrete}   \\
			\int g(x) f(x) \, dx & \text{if } X \text{ is continuous}
		\end{cases}
	\]
\end{definition}

Key properties:
\begin{itemize}
	\item $\E[aX + b] = a\E[X] + b$
	\item $\Var(X) = \E[(X - \E[X])^2] = \E[X^2] - \E[X]^2$
	\item $\Var(aX + b) = a^2 \Var(X)$
	\item $\Cov(X, Y) = \E[(X - \E[X])(Y - \E[Y])]$
\end{itemize}

\section{Multivariate Statistics}

For a $d$-dimensional random vector $\mathbf{X} = (X_1, \dots, X_d)$:

\begin{itemize}
	\item \textbf{Mean vector:} $\boldsymbol{\mu} = \E[\mathbf{X}]$
	\item \textbf{Covariance matrix:} $\Sigma_{ij} = \Cov(X_i, X_j)$
\end{itemize}

The multivariate normal distribution:

\[
	p(\mathbf{x}) = \frac{1}{(2\pi)^{d/2} |\Sigma|^{1/2}} \exp\left(-\frac12 (\mathbf{x} - \boldsymbol{\mu})^\top \Sigma^{-1} (\mathbf{x} - \boldsymbol{\mu})\right)
\]

This is the building block for Gaussian mixture models, linear discriminant analysis, and Gaussian processes.

\section{Bayes' Rule and Probabilistic Inference}

Bayes' rule is the foundation of probabilistic ML:

\[
	p(\theta \mid \mathcal{D}) = \frac{p(\mathcal{D} \mid \theta) \, p(\theta)}{p(\mathcal{D})}
\]

\begin{itemize}
	\item $p(\theta)$: \textbf{prior} --- what we believe about $\theta$ before seeing data.
	\item $p(\mathcal{D} \mid \theta)$: \textbf{likelihood} --- how probable the data is given $\theta$.
	\item $p(\theta \mid \mathcal{D})$: \textbf{posterior} --- updated belief after seeing data.
	\item $p(\mathcal{D})$: \textbf{evidence} --- marginal probability of the data.
\end{itemize}

\begin{example}[Coin Toss]
	Suppose we flip a coin 10 times and observe 7 heads. We model the coin's bias $\theta = P(\text{heads})$ with a Beta$(2,2)$ prior (peaked at 0.5). The posterior is:
	\[
		p(\theta \mid \text{7 heads}) \propto \text{Binomial}(7 \mid 10, \theta) \times \text{Beta}(\theta \mid 2, 2)
	\]
	\[
		= \text{Beta}(\theta \mid 2 + 7, 2 + 3) = \text{Beta}(\theta \mid 9, 5)
	\]
	The posterior mean is $\frac{9}{14} \approx 0.64$, shifted from the prior mean of 0.5 toward the observed 0.7.
\end{example}

\section{Maximum Likelihood Estimation}

The most common estimation principle: choose $\theta$ to maximize the likelihood of the data.

\[
	\hat{\theta}_{\text{MLE}} = \arg\max_\theta p(\mathcal{D} \mid \theta) = \arg\max_\theta \sum_{i=1}^n \log p(y_i \mid \mathbf{x}_i, \theta)
\]

\begin{example}[MLE for Linear Regression]
	Assume $y_i = \mathbf{w}^\top \mathbf{x}_i + \epsilon_i$, $\epsilon_i \sim \N(0, \sigma^2)$. The log-likelihood is:
	\[
		\log p(\mathcal{D} \mid \mathbf{w}) = -\frac{n}{2} \log(2\pi\sigma^2) - \frac{1}{2\sigma^2} \sum_{i=1}^n (y_i - \mathbf{w}^\top \mathbf{x}_i)^2
	\]
	Maximizing this is equivalent to minimizing $\sum_i (y_i - \mathbf{w}^\top \mathbf{x}_i)^2$ --- ordinary least squares.
\end{example}

\section{Maximum a Posteriori Estimation}

MAP estimation adds a prior:

\[
	\hat{\theta}_{\text{MAP}} = \arg\max_\theta \log p(\theta \mid \mathcal{D}) = \arg\max_\theta \big[\log p(\mathcal{D} \mid \theta) + \log p(\theta)\big]
\]

A Gaussian prior on $\mathbf{w}$ (with zero mean) is equivalent to L2 regularization:

\[
	\log p(\mathbf{w}) = -\frac{\lambda}{2} \|\mathbf{w}\|_2^2 + \text{const}
\]

Thus, ridge regression is the MAP estimate under a Gaussian prior.

\section{Hypothesis Testing and Confidence Intervals}

\subsection{Null Hypothesis Significance Testing}

When comparing two models A and B:

\begin{enumerate}
	\item \textbf{Null hypothesis} $H_0$: models A and B have equal performance.
	\item Compute a test statistic (e.g., difference in accuracy).
	\item Compute the $p$-value: probability of observing the statistic (or more extreme) under $H_0$.
	\item If $p < \alpha$ (typically 0.05), reject $H_0$.
\end{enumerate}

\begin{definition}[$p$-value]
	The \textbf{$p$-value} is the probability of obtaining results at least as extreme as those observed, assuming the null hypothesis is true. It is NOT the probability that $H_0$ is true.
\end{definition}

\subsection{Confidence Intervals}

A $95\%$ confidence interval for a parameter $\theta$ is an interval $[L, U]$ such that, if we repeated the experiment many times, $95\%$ of the intervals would contain the true $\theta$.

For model accuracy $\hat{a}$, an approximate confidence interval is:

\[
	\hat{a} \pm z_{\alpha/2} \sqrt{\frac{\hat{a}(1 - \hat{a})}{n}}
\]

where $z_{\alpha/2} = 1.96$ for $95\%$ confidence.

\section{Information-Theoretic Concepts}

\subsection{Entropy}

The entropy of a discrete random variable $X$ measures its uncertainty:

\[
	H(X) = -\sum_{x} p(x) \log p(x)
\]

\begin{itemize}
	\item $H(X) = 0$ if $X$ is deterministic.
	\item $H(X)$ is maximized for a uniform distribution.
\end{itemize}

\subsection{Kullback-Leibler Divergence}

The KL divergence measures how one distribution $p$ differs from another $q$:

\[
	D_{\text{KL}}(p \parallel q) = \sum_{x} p(x) \log \frac{p(x)}{q(x)}
\]

\begin{itemize}
	\item $D_{\text{KL}}(p \parallel q) \geq 0$, with equality iff $p = q$.
	\item Not symmetric: $D_{\text{KL}}(p \parallel q) \neq D_{\text{KL}}(q \parallel p)$.
	\item Used as the loss function in logistic regression and neural network classification (cross-entropy).
\end{itemize}

\section{Exercises}

\begin{exercise}
	Prove that $\Var(X) = \E[X^2] - \E[X]^2$.
\end{exercise}

\begin{exercise}
	You collect 100 data points from $\N(\mu, 1)$ and compute $\bar{x} = 3.2$. Compute a $95\%$ confidence interval for $\mu$ and interpret it. What would happen to the interval width if you collected 400 points?
\end{exercise}

\begin{exercise}
	Derive the MLE for the parameter $\lambda$ of a Poisson distribution: $p(k \mid \lambda) = \frac{\lambda^k e^{-\lambda}}{k!}$. Given data $\{k_1, \dots, k_n\}$, what is $\hat{\lambda}_{\text{MLE}}$?
\end{exercise}

\begin{exercise}
	Two classifiers achieve accuracies of $90\%$ and $92\%$ on a test set of 500 samples. Is the difference statistically significant at $\alpha = 0.05$? Compute the $p$-value using a normal approximation.
\end{exercise}
